\documentclass[12pt,letterpaper]{article}

\usepackage{amsmath}
\usepackage[russian]{babel}
\usepackage{graphicx} % для картинок
\usepackage{xcolor}

\title{Билет №5 \\ \small Множества и операции над ними. Булевы функции, КНФ, ДНФ. Базисы, теорема Поста. (ИТМО) \\ Алгебра логики. Булевы функции, канонические формы задания булевых функций. Понятие полной системы. Критерий полноты Поста. Минимизация булевых функций в классах нормальных форм. (СПбГУ) \\ Алгебра логики. Булевы функции, канонические формы задания булевых функций. Понятие полноты системы (СПбПУ)
}
\author{Сериков Владислав Александрович}
\date{30 марта 2026}

\begin{document}

\maketitle{}
\tableofcontents{}

\section{Теория множеств и операции над ними}

\textbf{\emph{Алгебра логики (алгебра высказываний)}} - раздел математической логики, в котором изучаются логические операции над высказываниями. \textbf{\emph{Математическая логика (теоретическая логика, символическая логика)}} - раздел математики, изучающий математические обозначения, формальные системы, доказуемость математических суждений, природу математического доказательства в целом, вычислимость и прочие аспекты оснований математики. \textbf{\emph{Высказывание}} в математической логике - предложение, выражающее суждение.

\textbf{\emph{Множество}} - это коллекция четко определенных и отличимых друг от друга объектов, которые называются элементами множества. Запись \(a \in A\) означает, что элемент \(a\) принадлежит множеству \(A\). Если элемент не принадлежит множеству, это обозначается как \(a \notin A\).

\textcolor{gray}{Примеры множеств:
\begin{itemize}
  \item Множество всех натуральных чисел: \(N = \{1, 2, 3, \ldots\} \)
  \item Множество всех целых чисел: \(Z = \{\ldots, -3, -2, -1, 0, 1, 2, 3, \ldots\}\).
  \item Множество всех действительных чисел: \(R\).
  \item Множество всех букв английского алфавита: \(A = \{a, b, c, \ldots, z\}\).
  \item Множество студентов группы: \(S = \{Ivanov, Petrov, Sidorov\}\).
\end{itemize}
Множества могут быть \textbf{\emph{конечными}} и \textbf{\emph{бесконечными}}.}

Если все элементы множества \(A\) принадлежат также множеству \(B\), то \(A\) называется \textbf{\emph{подмножеством}} \(B\), что обозначается \(A \subseteq B\). Множество \(A\) называется \textbf{\emph{строгим подмножеством}} множества \(B\), если каждый элемент множества \(A\) является элементом множества \(B\), однако не все элементы множества  \(B\)  содержатся в  \(A\), что обозначается \(A \subset B\). Если два множества содержат в точности одни и те же элементы, они считаются равными. Пустое множество  не содержит никаких элементов, обозначаемое символом \(\emptyset\). \textbf{\emph{Булево множество}} – множество состоящее из двух элементов, обозначаемое как \(B = {0, 1}\).

\subsection{Свойства подмножеств}

\begin{enumerate}
  \item \textbf{\emph{Рефлексивность:}} Любое множество является подмножеством самого себя, т.е. \(A \subseteq A\).
  \item \textbf{\emph{Антисимметричность:}} Если \(A \subseteq B\) и \(B \subseteq A\), то \(A = B\).
  \item \textbf{\emph{Транзитивность:}} Если \(A \subseteq B\) и \(B \subseteq C\), то \(A \subseteq C\).
  \item \textbf{\emph{Пустое множество:}} Пустое множество \(\emptyset\) является подмножеством любого множества: \(\emptyset \subseteq A\).
\end{enumerate}

\subsection{Операции над множествами}

\textbf{\emph{Операции над множествами}} - это способы получения новых множеств из уже существующих.

\subsubsection{Объединение}

\textbf{\emph{Объединением двух множеств \(A\) и \(B\)}}, обозначаемым \(A \cup B\), называется множество, содержащее все элементы, которые принадлежат хотя бы одному из множеств \(A\) или \(B\).
\[A \cup B = x \vert x \in A \vee x \in B\]
Эта операция \textbf{\emph{коммутативна}}: \(A \cup B = B \cup A\), и \textbf{\emph{ассоциативна}}: \((A \cup B) \cup C = A \cup (B \cup C)\).

\subsubsection{Пересечение}

\textbf{\emph{Пересечением двух множеств \(A\) и \(B\)}}, обозначаемым \(A \cap B\), называется множество, содержащее все элементы, которые принадлежат обоим множествам одновременно.
\[A \cap B = x \vert x \in A \wedge x \in B\]
Эта операция \textbf{\emph{коммутативна}}, \textbf{\emph{ассоциативна}} и \textbf{\emph{дистрибутивна}} относительно объединения и наоборот.
\[A \cap (B \cup C) = (A \cap B) \cup (A \cap C)\]
\[A \cup (B \cap C) = (A \cup B) \cap (A \cup C)\]

\subsubsection{Разность}

\textbf{\emph{Разностью двух множеств \(A\) и \(B\)}}, обозначаемой \(A \setminus B\), называется множество, содержащее все элементы, которые принадлежат \(A\), но не принадлежат \(B\).
\[A \setminus B = x \vert x \in A \wedge x \notin B\]
Эта операция не является \textbf{\emph{ни коммутативной}}, \textbf{\emph{ни ассоциативной}}.

\subsubsection{Симметрическая разность}

\textbf{\emph{Симметрическая разность}}, обозначаемая \(A \triangle B\), представляет собой объединение разностей \(A \setminus B\) и \(B \setminus A\). Элемент принадлежит симметрической разности тогда и только тогда, когда он принадлежит в точности одному из двух множеств.
\[A \triangle B = (A \setminus B) \cup (B \setminus A)\]

\subsubsection{Дополнение}

\textbf{\emph{Дополнением множества A относительно некоторого универсального множества U}}, которое содержит все рассматриваемые элементы, называется множество \(A^c = U \setminus A\). Элемент принадлежит дополнению \(A\), если он не принадлежит \(A\).

\subsection{Законы де Моргана}

\textcolor{gray}{Они описывают взаимосвязь между объединением, пересечением и операцией дополнения.}
Для любых двух множеств \(A\) и \(B\) справедливы следующие тождества:
\begin{enumerate}
  \item \((A \cup B) ^ c = A ^ c \cap B ^ c\).
  \item \((A \cap B) ^ c = A ^ c \cup B ^ c\).
\end{enumerate}
\textcolor{gray}{Докажем первое тождество. Возьмем произвольный элемент \(x \vert x \in  (A \cup B) ^ c \Leftrightarrow x \notin (A \cup B)\) (по определению дополнения) \(\Leftrightarrow \neg (x \in (A \cup B)) \Leftrightarrow \neg (x \in A \vee x \in B)\) (по определению объединения). \(\Leftrightarrow (x \notin A) \wedge (x \notin B) \Leftrightarrow (x \in A ^ c) \wedge (x \in B ^ c)\) (по определению дополнения). \(\Leftrightarrow x \in (A ^ c \cap B ^ c)\) (по определению пересечения). Поскольку \(x\) был произвольным, это доказывает, что множества \((A \cup B) ^ c\)  и \(A ^ c \cap B ^ c\)  содержат одни и те же элементы, а значит, они равны.
Второе тождество доказывается аналогично.}

\section{Булевы функции}

\textbf{\emph{Булева алгебра}} состоит из непустого множества \(B\), содержащего, как минимум, два различных элемента, обычно обозначаемых как 0 и 1, на котором определены три основные операции: бинарная операция \textbf{\emph{конъюнкции (\(\wedge\))}}, бинарная операция \textbf{\emph{дизъюнкции (\(\vee\))}} и унарная операция \textbf{\emph{отрицания (\(\neg\))}}. Наиболее важным примером булевой алгебры в дискретной математике является двухэлементная булева алгебра, где \(B = \{0, 1\}\).

\textbf{\emph{Булева функция от \(n\) переменных}} - это функция \(f : B ^ n \rightarrow B\), то есть правило, которое каждой совокупности из \(n\) значений переменных \(x_1, x_2, \ldots, x_n\)  из множества \{0,1\} ставит в соответствие единственное значение \(f(x_1, x_2, \ldots, x_n)\) из того же множества \{0,1\}.

\textbf{\emph{Арность (arity) функции}} - количество ее аргументов. Каждая булева функция арности \(n\) полностью определяется заданием своих значений на \(dom(f)\), то есть на всех булевых векторах длины \(n\). Число таких векторов равно \(2 ^ n\). Поскольку на каждом векторе булева функция может принимать значение либо 0, либо 1, то количество всех \(n\)-арных булевых функций равно \(2 ^ {2 ^ n}\). Обозначается \(P_2 (n)\).

\begin{itemize}
  \item \textbf{\emph{Нульарные функции}} - это булевые функции, которые не зависят ни от каких переменных. Существует всего две нульарные функции: \(f() = 0\) и \(g() = 1\). В логике нульарные функции соответствуют булевым константам.
  \item \textbf{\emph{Унарные функции}} - это булевые функции, определённые на одной переменной. Существует четыре возможных унарных функции:
  \begin{itemize}
    \item Тождественная функция \(f(x) = x\), которая просто возвращает значение переменной.
    \item Отрицание \(g(x) = \neg x\), которое обращает значение переменной.
    \item Константа 0: \(h(x) = 0\), которая всегда возвращает 0, независимо от значения переменной.
    \item Константа 1: \(k(x) = 1\), которая всегда возвращает 1, независимо от значения переменной.
  \end{itemize}
  \item \textbf{\emph{Бинарные функции}} - это булевые функции от двух переменных.
  \item \textbf{\emph{Тернарные функции}} - это булевые функции от трех переменных.
\end{itemize}

\subsection{Основные операции над булевыми функциями}

\textbf{\emph{Таблица истинности}} - это таблица, в которой перечислены все \(2 ^ n\)  возможных наборов значений аргументов \(x_1, x_2, \ldots, x_n\)  и соответствующие значения функции \(f\) на этих наборах. Любая булева функция однозначно определяется своей таблицей истинности, и, наоборот, любая таблица истинности размером \(2 ^ n\)  задает некоторую булеву функцию от \(n\) переменных.

\begin{center}
\fontsize{10pt}{10pt}
\begin{tabular}{|c|c|c|c|c|c|c|c|c|c|c|c|c|c|}
  \hline
  \(x\) & \(y\) & \(x \wedge y\) & \(x \vee y\) & \(\neg x\) & \(x \rightarrow y\) & \(x \oplus y\) & \(x \leftrightarrow y\) & \(x \uparrow y\) & \(x \downarrow y\) & \(x \leftarrow y\) & \(I(x)\) & \(0(x)\) & \(1(x)\) \\
  \hline
  0 & 0 & 0 & 0 & 1 & 1 & 0 & 1 & 1 & 1 & 1 & 0 & 0 & 1 \\
  0 & 1 & 0 & 1 & 1 & 1 & 1 & 0 & 1 & 0 & 0 & 0 & 0 & 1 \\
  1 & 0 & 0 & 1 & 0 & 0 & 1 & 0 & 1 & 0 & 1 & 1 & 0 & 1 \\
  1 & 1 & 1 & 1 & 0 & 1 & 0 & 1 & 0 & 0 & 1 & 1 & 0 & 1 \\
  \hline
\end{tabular}
\end{center}

Для \(\forall x \in B\) его булевой степенью функции называется:
\begin{equation*}
x ^ \sigma =
  \begin{cases}
    x, & \mbox{если } \sigma = 1 \\
    \neg x, & \mbox{если } \sigma = 0
  \end{cases}
\end{equation*}
Легко видеть, что \(x ^ \sigma  = x \oplus \sigma  \oplus 1 = x ^ \sigma  \vee x ^ \sigma\). Для любого \(x = (x_1, x_2, \ldots, x_n)\) и любого \(\sigma = (\sigma_1,  \sigma_2,  \ldots, \sigma_n)\) произведение:
\[k_\sigma(x) = x_1 ^ {\sigma_1} \cdots x_n ^ {\sigma_n}\]
булевых степеней \(\sigma_1, \ldots, \sigma_n\) переменных \(x_1, \ldots, x_n\) называется \textbf{\emph{элементарной конъюнкцией}}, ассоциированной с булевым набором \(\sigma = \sigma_1, \sigma_2, \ldots, \sigma_n\) . Число переменных, входящих в конъюнкцию \(k_\sigma(x)\) , называется \textbf{\emph{рангом этой конъюнкции}}. Элементарная конъюнкция \(x_1 ^ {\sigma_1}, x_2 ^ {\sigma_2}, \ldots, x_n ^ {\sigma_n}\) обращается в 1 тогда и только тогда, когда одновременно выполняются равенства \(x_1 = \sigma_1, x_2 = \sigma_2, \ldots, x_n = \sigma_n\). Для функции \(k_\sigma(x)\) справедливо соотношение:
\begin{equation*}
k_\sigma(x) = 
  \begin{cases}
    1, \mbox{если } \sigma = 1 \\
    0, \mbox{если } \sigma = 0
  \end{cases}
\end{equation*}

Отсюда немедленно следует, что \(k_a(x) \cdot k_b(x) = 0\) при \(a \neq b\).

\subsection{Теорема о разложении функции по переменным}

Пусть даны функция \(f(x_1, \ldots , x_n)\) и число \(k, 1 \leq k \leq n\). Тогда функцию \(f\) можно представить в следущей форме:
\[f(x_1, \ldots, x_n) =  \bigvee_{(\sigma_1, \ldots, \sigma_n)} x_1 ^ {\sigma_1} \cdots x_n ^ {\sigma_n} f(\sigma_1, \ldots, \sigma_k, x_{k+1}, \ldots, x_n), \]
где дизъюнкция берется по всем наборам \((\sigma_1, \ldots, \sigma_n)\) значений переменных \((x_1, \ldots, x_n)\).

Доказательство: Из определения элементарной конъюнкции легко следует, что:
\[\bigvee_{\sigma_1, \ldots, \sigma_k} \alpha_1 ^ {\sigma_1} \cdots \alpha_k ^ {\sigma_k} \cdot f(\sigma_1, \ldots, \sigma_k, \alpha_{k+1}, \ldots, \alpha_n) = \alpha_1 ^ {\alpha_1} \cdots \alpha_k ^ {\alpha_k} \cdot f(\alpha_1, \ldots, \alpha_n) = f(\alpha_1, \ldots, \alpha_n)\]

\textcolor{red}{РАЗОБРАТЬ}

\subsection{Нормальные формы}

\textbf{\emph{Литерал}} - это булева переменная или ее отрицание. Например, для переменной \(x\), литералы - \(x\) и \(\neg x\).

\textbf{\emph{Минтерм}} - это конъюнкция всех переменных функции или их отрицаний. Минтерм принимает значение 1 для одного и только одного набора значений переменных. Минтерм зависит от всех переменных булевой функции и имеет вид:
\[M = x_1 ^ {\alpha_1} \wedge x_2 ^ {\alpha_2} \wedge \ldots \wedge x_n ^ {\alpha_n}\]
где \(\alpha_i \in \{0, 1\}\), а \(x_i ^ 0 = \neg x_i, x_i ^ 1 = x_i\). Минтермы используются для построения дизъюнктивной нормальной формы (ДНФ).

\textbf{\emph{Макстерм}} - это дизъюнкция всех переменных или их отрицаний. Макстерм принимает значение 0 для одного и только одного набора значений переменных и имеет вид:
\[M' = x_1 ^ {\alpha_1} \vee x_2 ^ {\alpha_2} \vee \ldots \vee x_n ^ {\alpha_n}\]
где \(\alpha_i \in \{0, 1\}\), а \(x_i ^ 0 = \neg x, x_i ^ 1 = x_i\). Минтермы используются для построения конъюктивной нормальной формы (КНФ).

Для систематического анализа и упрощения булевых функций используются различный формы и записи, известные как \textbf{\emph{нормальные формы}}. ключевыми среди них являются \textbf{\emph{дизъюнктивная нормальная форма (ДНФ)}} и \textbf{\emph{конъюктивная нормальная форма (КНФ)}}.
\begin{itemize}
  \item \textbf{\emph{Дизъюнктивная нормальнаяя форма (ДНФ)}} - это выражение булевой функции в виде дизъюнкции нескольких конъюнкций литералов (переменных и их отрицаний). ДНФ может иметь вид:
  \[f(x_1, x_2, \ldots, x_n) = (l_{11} \wedge l_{12} \wedge \ldots \wedge l_{1n}) \vee \ldots \vee (l_{m1} \wedge l_{m2} \wedge \ldots \wedge l_{mn}),\]
  где, \(l_{ij}\) - литералы.
  \item \textbf{\emph{Конъюктивная нормальная форма}} - выражение булевой функции в виде конъюнкции нескольких дизъюнкций литералов. КНФ может иметь вид:
  \[f(x_1, x_2, \ldots, x_n) = (l_{11} \vee l_{12} \vee \ldots \vee l_{1n}) \wedge \ldots \wedge (l_{m1} \vee l_{m2} \vee \ldots \vee l_{mn}),\]
  где, \(l_{ij}\) - литералы.
\end{itemize}

Любая булева функция, не являющаяся тождественным нулем, может быть представлена в виде ДНФ. Любая булева функция, не являющаяся тождественной единицей, может быть представлена в виде КНФ.

Чтобы получить уникальное и полное представление любой булевой функции, вводятся канонические формы. К ним относятся \textbf{\emph{совершенная дизъюнктивная нормальная форма (СДНФ)}} и \textbf{\emph{совершенная конъюнктивная нормальная форма (СКНФ)}}.

\textbf{\emph{Совершенная дизъюнктивная нормальная форма (СДНФ)}} - это ДНФ, в которой нет двух одинаковых конъюнктов и все переменные исходного выражения присутствуют в каждом конъюнкте. СДНФ представляет функцию как дизъюнкцию всех минтермов, на которых функция принимает значение 1:
\[f(x_1, \ldots, x_n) = \bigvee_{(\sigma_1, \ldots, \sigma_n):f(\sigma_1, \ldots, \sigma_n) = 0} (x_1 ^ {\sigma_1} \wedge \ldots \wedge x_n ^ {\sigma_n})\]

Алгоритм построения СДНФ по талице истинности следующий:

\begin{enumerate}
  \item Выбрать все строки таблицы истинности, в которых значение функции \(f\) равно 1.
  \item Для каждой такой строки составить минтерм. Если переменная в этой строке имеет значение 1, она входит в минтерм без отрицания; если 0 - с отрицанием.
  \item Составить дизъюнкцию всех полученных минтермов.
\end{enumerate}

Пример: пусть дана функция \(f(x, y, z)\) со следующей таблицей истинности

\begin{center}
\begin{tabular}{|c|c|c|c|}
  \hline
  \(x\) & \(y\) & \(z\) & \(f(x, y, z)\) \\
  \hline
  0 & 0 & 0 & 0 \\
  \hline
  0 & 0 & 1 & 1 \\
  \hline
  0 & 1 & 0 & 0 \\
  \hline
  0 & 1 & 1 & 1 \\
  \hline
  1 & 0 & 0 & 1 \\
  \hline
  1 & 0 & 1 & 0 \\
  \hline
  1 & 1 & 0 & 0 \\
  \hline
  1 & 1 & 1 & 1 \\
  \hline
\end{tabular}
\end{center}

Функция равна 1 в строках 1, 3, 4 и 7 (нумерация с нуля). СДНФ будет дизъюнкцей минтермов для этих строк:

\begin{itemize}
  \item Строка 1 (001): \((\neg x \wedge \neg y \wedge z)\)
  \item Строка 3 (011): \((\neg x \wedge y \wedge z)\)
  \item Строка 4 (100): \((x \wedge \neg y \wedge \neg z)\)
  \item Строка 7 (111): \((x \wedge y \wedge z)\)
\end{itemize}

СДНФ: \(f(x, y, z) = (\neg x \wedge \neg y \wedge z) \vee (\neg x \wedge y \wedge z) \vee (x \wedge \neg y \wedge \neg z) \vee (x \wedge y \wedge z)\)

\textbf{\emph{Совершенная конъюктивная нормальная форма (СКНФ)}} - это КНФ, в которой нет двух одинаковых дизъюнктов и все переменные присутствуют в каждом дизъюнкте. СКНФ представляет функцию как конъюнкцию всех макстермов, на которых функция ринимает значение 0:

\[f(x_1, x_2, \ldots, x_n) = \bigwedge_{(\sigma_1, \ldots, \sigma_n):f(\sigma_1, \ldots, \sigma_n) = 0} (x_1 ^ {\sigma_1} \vee \ldots \vee x_n ^ {\sigma_n})\]

Алгоритм построения СКНФ по таблице истинности:

\begin{enumerate}
  \item Выбрать все строки таблицы истинности, в которых значение функции \(f\) равно 0.
  \item Для каждой такой строки составить макстерм. Если переменная в этой строке имеет значение 0, она входит в макстерм без отрицания; если 1 - с отрицанием.
  \item Составить конъюнкцию всех полученных макстермов.
\end{enumerate}

Применяя этот алгоритм к нашей функции \(f(x, y, z)\), которая равна 0 в строках 0, 2, 5, 6, получаем СКНФ:

\begin{itemize}
  \item Строка 0 (000): \((x \vee y \vee z)\)
  \item Строка 2 (010): \((x \vee \neg y \vee z)\)
  \item Строка 5 (101): \((\neg x \vee y \vee \neg z)\)
  \item Строка 6 (110): \((\neg x \vee \neg y \vee z)\)
\end{itemize}

СКНФ: \(f(x, y, z) = (x \vee y \vee z) \wedge (x \vee \neg y \vee z) \wedge (\neg x \vee y \vee \neg z) \wedge (\neg x \vee \neg y \vee z)\)

Существуют другие разложения. Например, полином Жегалкина. Булева функция может быть представлена в виде Полинома Жегалкина, который является уникальной нормальной формой булевой функции. Полином Жегалкина - это выражение, состоящее из суммы (операция XOR) произведений булевых переменных и констант 0 или 1 Важным свойством полинома Жегалкина является то, что он не содержит операции отрицания.

Полином Жегалкина для булевой функции \(f(x_1, x_2, \ldots, x_n)\) может быть записан следующим образом:

\[f(x_1, x_2, \ldots, x_n) = c_0 \oplus c_1 x_1 \oplus c_2 x_2 \oplus \ldots \oplus c_n x_n \oplus c_{12} x_1 x_2 \oplus \ldots \oplus c_{12 \ldots n} x_1 x_2 \ldots x_n\]

Каждое произведение в этом полиноме соответствует одному подмножеству переменных, где результат произведения истинен, если все переменные в этом подмножестве равны 1. Например, для функции трёх переменных, \(f(x_1, x_2, x_3) = x_1 \wedge \neg x_2 \wedge x_3\) полином Жегалкина будет содержать моном \(x_1 x_3\) (так как \(x_2\) равно 0, оно не входит в моном).

Свойства полинома Жегалкина:

\begin{enumerate}
  \item \textbf{\emph{Линейность:}} полином Жегалкина является линейным по каждой переменной, так как операция сложения реализуется с помощью операции XOR.
  \item \textbf{\emph{Уникальность:}} для каждой булевой функции существует единственный соответствующий ей полином Жегалкина. 
\end{enumerate}

\section{Замыкание, полнота, Критерий Поста}

Замыкание показывает, какие новые функции можно получить из данного набора функций, но не объясняет, как именно это делается. Суперпозиция описывает механизм создания этих новых функций путём подстановки. Полнота систем непосредственно связана с возможностью построения любой функции через суперпозиции, поэтому важно сначала объяснить сам процесс суперпозиции, а затем перейти к обсуждению полноты, где это понятие уже используется.

\subsection{Суперпозиция}

\textbf{\emph{Суперпозиция}} - это процесс построения новой булевой функции путём подстановки одной или нескольких булевых функций в другую булеву функцию. Если даны функции \(f(x_1, x_2, \ldots, x_n)\) и \(g_1(x_1, x_2, \ldots, x_n)\), \(g_2(x_1, x_2, \ldots, x_n)\), \(\ldots\), \(g_k(x_1, x_2, \ldots, x_n)\), то суперпозиция функций \(f\) и \(g_1\), \(g_2\), \(\ldots\), \(g_n\) записывается как:

\[h(x_1, x_2, \ldots, x_n) = f(g_1(x_1, x_2, \ldots, x_n), g_1(x_1, x_2, \ldots, x_n), \ldots, g_k(x_1, x_2, \ldots, x_n))\]

Таким образом, суперпозиция строит новую функцию, используя результаты других функций в качестве входных данных для первой функции.

\subsubsection{Простая и сложная суперпозиции}

\begin{itemize}
  \item \textbf{\emph{Простая суперпозиция}} - это суперпозиция, в которой каждая подставляемая функция зависит только от одной переменной. Например, функция \(f(g_1(x_1), g_2(x_2))\) является простой суперпозицией.
  \item \textbf{\emph{Сложная суперпозиция}} - это суперпозиция, в которой подставляемые функции могут зависеть от нескольких переменных одновременно. Например, функция \(f(g_1(x_1, x_2), g_2(x_1, x_2))\) является сложной суперпозицией. 
\end{itemize}

\subsubsection{Свойства суперпозиций булевых функций}

\begin{itemize}
  \item \textbf{\emph{Ассоциативность:}} Суперпозиция булевых функций является ассоциативной. Это значит, что порядок выполнения операций суперпозиции не влияет на конечный результат:
  \[f(g(h(x_1, \ldots, x_n))) = (f \circ g) \circ h = f \circ (g \circ h)\]
  \item \textbf{\emph{Коммутативность для некоторых операций:}} Суперпозиция коммутативна, если используемые в ней операции являются коммутативными (например, дизъюнкция или конъюнкция).
\end{itemize}

\subsection{Замыкание множества булевых функций}

Пусть \(F\) - некоторое множество булевых функций. Замыканием множества \(F\), обозначаемым через  \(\left[ F \right]\), называется наименьшее множество функций, которое содержит \(F\) и замкнуто относительно всех операций суперпозиции (композиции функций). Множество булевых функций R называется (функционально) замкнутым множеством, если оно совпадает со своим замыканием, т. е. \(F = \left[ F \right]\).
Формально:
\[\left[ F \right] = \{f \vert f \mbox{ может быть представлена как суперпозиция функций из } F\}\].

\subsubsection{Замкнутые классы}

\begin{enumerate}
  \item Функции, сохраняющие ноль (\(T_0\)): \(f(0, \ldots, 0) = 0\).

  Доказательство:

  \[f(0, \ldots, 0) = f_0(f_1(0, \ldots, 0), \ldots, f_k(0, \ldots, 0)) = f_0(0, \ldots, 0) = 0\]

  Реализуемая формулой \(f_0(f_1, \ldots, f_k)\) функция \(f\) также сохраняет нуль, следовательно \(T_0\) замкнуто.

  \item Функции, сохраняющие единицу (\(T_1\)): \(f(1, \ldots, 1) = 1\).
  \item Линейные (\(L\)).
  \item Самодвойственные (\(S\)).
  \item Монотонные (\(M\)).
\end{enumerate}

\subsubsection{Линейные функции}

Булева функция \(f(x_1, x_2, \ldots, x_n)\) называется линейной, если она может быть записана в виде линейного полинома:
\[f(x_1, x_2, \ldots, x_n) = a_0 \oplus a_1 x_1 \oplus a_2 x_2 \oplus \ldots \oplus a_n x_n,\]
где \(a_0, a_1, a_2, \ldots, a_n \in \{0, 1\}\), а \(\oplus\) - операция сложения по модулю 2 (исключающее ИЛИ, XOR).

\textbf{Свойства линейных функций}

\begin{enumerate}
  \item Линейные функции сохраняют сложение по модулю 2.
  \item Композиция двух линейных функций также является линейной.
\end{enumerate}

\subsubsection{Двойственные функции}

Булева функция \(f(x_1, x_2, \ldots, x_n)\) называется двойственной функции \(f'(x_1, x_2, \ldots, x_n)\), если выполняется следующее условие:

\[f(x_1, x_2, \ldots, x_n) = \neg f'(\neg x_1, \neg x_2, \ldots, \neg x_3)\]

\subsubsection{Самодвойственные функции}

Функция \(f(x_1, x_2, \ldots, x_n)\) называется самодвойственной, если она совпадает со своей двойственной функцией, то есть:

\[f(x_1, x_2, \ldots, x_n) = \neg f(\neg x_1, \neg x_2, \ldots, \neg x_n)\]

\textbf{\emph{Докажем замкнутость множества самодвойственных функций.}} Пусть \(f_0\), \(f_1\), \(\ldots\), \(f_k\) - произвольные самодвойственные функции. Рассмотрим новую функцию \(f = f_0(f_1, \ldots, f_k)\). Так как добавление фиктивной переменной оставляет самодвойственную функцию самодвойственной, то без ограничения общности будем полагать, что все функции \(f_i\) зависят от одних и тех же переменных \(x_1, \ldots, x_n\). Тогда \(f(\neg x_1, \ldots, \neg x_n) = f_0(f_1(\neg x_1, \ldots, \neg x_n), \ldots, f_k(\neg x_1, \ldots, \neg x_n)) == f_0(\neg f_1(x_1, \ldots, x_n), \ldots, \neg f_k(x_1, \ldots, x_n)) == \neg f_0(f_1(x_1, \ldots, x_n), \ldots, f_k(x_1, \ldots, x_n)) = \neg f(x_1, \ldots, x_n)\). Следовательно, функция \(f\) — самодвойственная. Таким образом, множество \(S\) замкнуто.

\subsubsection{Монотонные функции}

Булева функция \(f(x_1, x_2, \ldots, x_n)\) называется монотонной, если для любых двух наборов значений переменных \((x_1, x_2, \ldots, x_n)\) и \((y_1, y_2, \ldots, y_n)\), таких что \(x_i \leq y_i\) для всех \(i\), выполняется неравенство:

\[f(x_1, x_2, \ldots, x_n) \leq f(y_1, y_2, \ldots, y_n)\]

\textbf{\emph{Докажем замкнутость множества монотонных функций.}} Пусть \(f_0, \ldots, f_k\) - произвольные монотонные функции. Очевидно, что добавление фиктивной переменной оставляет монотонную функцию монотонной. Поэтому без ограничения общности будем полагать, что все функции \(f_i\) зависят от одних и тех же переменных \((x_1, \ldots, x_n)\). Пусть \(a, b\) - такие наборы из \(B^n\), что \(a \leq b\). Рассмотрим новую функцию \(f = f_0(f_1, \ldots, f_k)\). Так как \(f_i(a) \leq f_i(b)\), то \((f_1(a), \ldots, f_k(a)) \leq (f_1(b), \ldots, f_k(b))\), и поэтому \(f(a) = f_0(f_1(a), \ldots, f_k(a)) \leq f_0(f_1(b), \ldots, f_k(b)) = f(b)\). Следовательно, функция \(f\) - монотонная. Таким образом, множество \(M\) замкнуто.

\textcolor{red}{РАЗОБРАТЬ}

\subsection{Полнота систем}

\end{document}